\documentclass[book.tex]{subfiles}

\begin{document}
\chapter{The Bloch Sphere}

\label{chap:bloch_sphere}
\noindent\rule{11cm}{0.4pt} \\
\textbf{Prerequisites:} \autoref{chap:quantum_mechanics} \\
\textbf{Difficulty Level:} ** \\
\noindent\rule{11cm}{0.4pt}
\vspace{5mm}

The Bloch sphere is a representation of the state space of a two level quantum system. We can write such a system by the equation

\begin{align}
    |\psi \rangle &= \alpha | 0 \rangle + \beta | 1 \rangle \\
\end{align}

We know that the norm of this state has to be $1$, so

\begin{align}
    \langle \psi | \psi \rangle &= \alpha^2 + \beta^2 = 1
\end{align}

Note also that the coefficients $\alpha$ and $\beta$ are only defined up to a complex phase factor $e^{i\phi}$ if
\begin{align}
    \alpha' &= e^{i\phi} \alpha \\
    \beta' &= e^{i \phi} \beta \\
    |\psi' \rangle &= \alpha' |0 \rangle + \beta' | 1 \rangle \\
\end{align}
Then both $|\psi \rangle$ and $|\psi' \rangle$ represent the same quantum state. We can then chose to multiply by a phase factor so that the coefficient of $|0 \rangle$ is real and nonnegative. This allows us to write

\begin{align}
    | \psi \rangle &= \cos \theta | 0 \rangle + e^{i \phi} \sin \theta | 1 \rangle 
\end{align}
Here $\theta$ must range between $0$ and $\pi/2$ for both $\cos \theta$ and $\sin \theta$ to be nonnegative. $\phi$ ranges between $0$ and $2 \pi$
We can go one step further and do a change of variables to replace $\theta$ by $\theta/2$. Then we have the following expression
\begin{align}
  | \psi \rangle &= \cos (\theta/2) | 0 \rangle + e^{i \phi} \sin (\theta/2) | 1 \rangle 
\end{align}
where $\theta$ ranges between $0$ and $\pi$ and $\phi$ ranges from $0$ to $2 \pi$. These ranges allow us to interpret $\theta$ and $\phi$ as spherical coordinates! We can use them to write a unique point on the unit sphere as 
\begin{align}
    (\sin \theta \cos \phi, \sin \theta \sin \phi, \cos \theta)
\end{align}
This spherical representation of $\psi$ is denoted as the Bloch sphere as the diagram below shows.
\begin{marginfigure}
    \centering
    % \includegraphics{figures/Physics/quantum_computing/bloch_sphere/bloch_sphere.png}
    \includegraphics{figures/Physics/quantum_computing/bloch_sphere/bloch_sphere1.png}
    \caption{The Bloch shere provides a representation of the state space of a two level quantum system.} %TODO: This is from \url{https://en.wikipedia.org/wiki/Bloch_sphere}. Replace with our own figure.}
    \label{fig:my_label}
\end{marginfigure}

To gain intuition for the Bloch sphere, it is useful to plot a few representative states. Where do the basis states $|0 \rangle$ and $| 1 \rangle$ land for example? We start by rewriting both of these states in the canonical representation above
\begin{align}
    |0 \rangle &= \cos (0/2) | 0 \rangle + e^{i \phi} \sin(0/2) | 1 \rangle \\
    |1 \rangle &= \cos(\pi/2) | 0 \rangle + e^{i \phi} \sin(\pi/2) | 1 \rangle 
\end{align}
Here $\phi$ can be an arbitrary value between $0$ and $2\pi$. We see then that $|0 \rangle$ lands at the North pole and $|1 \rangle$ at the South pole of the Bloch sphere.

\end{document}